\documentclass[aspectratio=169,10pt]{beamer}
\usepackage[utf8]{inputenc}
\usepackage[T1]{fontenc}
\usepackage[spanish]{babel}
\usepackage{lmodern,amsmath,amssymb,tikz,booktabs,array}
\usetikzlibrary{arrows.meta,calc,positioning,decorations.pathreplacing}
\definecolor{azul}{RGB}{14,72,130}\definecolor{verde}{RGB}{0,130,100}\definecolor{morado}{RGB}{115,50,170}\definecolor{naranja}{RGB}{225,110,25}
\usetheme{Madrid}\setbeamertemplate{headline}{}\setbeamertemplate{blocks}[rounded][shadow=false]\setbeamerfont{frametitle}{size=\large,series=\bfseries}\setbeamercolor{background canvas}{bg=white}\setbeamercolor{normal text}{fg=black!85}\setbeamersize{text margin left=11mm,text margin right=11mm}\setbeamercolor{structure}{fg=azul}\setbeamercolor{frametitle}{bg=azul,fg=white}\setbeamercolor{block title}{bg=azul,fg=white}\setbeamercolor{block body}{bg=blue!5,fg=black}\setbeamercolor{block title example}{bg=verde,fg=white}\setbeamercolor{block body example}{bg=green!5,fg=black}
\setbeamertemplate{navigation symbols}{}\setbeamertemplate{footline}[frame number]
\newcommand{\res}[1]{\[\boxed{\color{morado}#1}\]}\newcommand{\etapa}[1]{\par\medskip{\color{azul}\bfseries #1}\par}
\newcommand{\datos}[1]{\textbf{Datos:}\quad #1\par\medskip}
\newcommand{\paso}[2]{\textbf{#1}\[ #2 \]}
\newcommand{\auto}[3]{\begin{tikzpicture}[scale=.8,>=Stealth]\draw[thick,gray](-.5,0)--(7,0);\draw[fill=blue!15,draw=azul,rounded corners=2pt](#1,.15) rectangle ++(1.6,.65);\draw[fill=black](#1+.35,.12)circle(.12);\draw[fill=black](#1+1.25,.12)circle(.12);\draw[->,very thick,verde](#1+.8,1.1)--++(#2,0) node[above,midway]{#3};\end{tikzpicture}}
\title{Movimiento rectilíneo uniformemente acelerado y movimiento vertical}
\subtitle{MRUA, caída libre y lanzamiento vertical}
\author{Clases de Matemática Paraíso}\date{Física -- Educación Diversificada}
\begin{document}
\begin{frame}\titlepage\end{frame}
\begin{frame}{Objetivos de aprendizaje}\begin{enumerate}\item Reconocer las características del MRU y del MRUA.\item Diferenciar velocidad, aceleración y desplazamiento.\item Resolver problemas de aceleración y frenado horizontal.\item Resolver problemas de caída libre y lanzamientos verticales.\item Interpretar el signo de la velocidad y de la aceleración.\end{enumerate}\end{frame}
\section{MRUA horizontal}
\begin{frame}{Recordatorio: MRU y MRUA}\begin{block}{Movimiento rectilíneo uniforme (MRU)}La velocidad es constante y la aceleración es cero.\vspace{2mm}\[v=\frac{d}{t}\qquad d=vt\]\end{block}\begin{block}{Movimiento rectilíneo uniformemente acelerado (MRUA)}La aceleración es constante: la velocidad cambia la misma cantidad en intervalos de tiempo iguales.\vspace{2mm}\[a=\frac{v-v_0}{t}\]\end{block}\end{frame}
\begin{frame}{Variables y unidades}\centering\renewcommand{\arraystretch}{1.35}\begin{tabular}{lll}\toprule Magnitud& Símbolo& Unidad SI\\\midrule Desplazamiento&$d$&$\mathrm{m}$\\Velocidad inicial&$v_0$&$\mathrm{m/s}$\\Velocidad final&$v$&$\mathrm{m/s}$\\Aceleración&$a$&$\mathrm{m/s^2}$\\Tiempo&$t$&$\mathrm{s}$\\\bottomrule\end{tabular}\par\bigskip\alert{Una velocidad o un desplazamiento puede ser negativo según el eje elegido.}\end{frame}
\begin{frame}{Formulario de MRUA}\begin{columns}[T]\column{.5\textwidth}\begin{block}{Con el tiempo}\[v=v_0+at\]\[d=v_0t+\frac{at^2}{2}\]\[d=\left(\frac{v_0+v}{2}\right)t\]\end{block}\column{.5\textwidth}\begin{block}{Sin tiempo}\[v^2=v_0^2+2ad\]\end{block}\begin{block}{Definición de aceleración}\[a=\frac{v-v_0}{t}\]\end{block}\end{columns}\smallskip\centering\small Las ecuaciones suponen aceleración constante y $d$ es desplazamiento con signo.\end{frame}
\begin{frame}{Convención de signos: horizontal}\begin{center}\begin{tikzpicture}[>=Stealth]\draw[->,very thick,azul](-4,0)--(4,0)node[right]{Este $(+)$};\draw[->,very thick,naranja](0,-.8)--(-3,-.8)node[left]{Oeste $(-)$};\node[above] at (0,0){Origen};\end{tikzpicture}\end{center}\begin{itemize}\item Hacia la derecha: velocidad positiva.\item Hacia la izquierda: velocidad negativa.\item La aceleración positiva apunta a la derecha; la negativa, a la izquierda.\item \textbf{Frenar} significa que aceleración y velocidad tienen sentidos opuestos, no necesariamente que $a<0$.\end{itemize}\end{frame}
\begin{frame}{Situación 1: vehículo que parte del reposo}\centering\auto{1}{3}{$a=2,5\,\mathrm{m/s^2}$}\par\medskip\begin{block}{Enunciado}Un automóvil parte del reposo y acelera uniformemente a $2,5\,\mathrm{m/s^2}$ durante $8\,\mathrm{s}$. Determine la velocidad final y el desplazamiento.\end{block}\datos{$v_0=0\,\mathrm{m/s}$; $a=2,5\,\mathrm{m/s^2}$; $t=8\,\mathrm{s}$.}\end{frame}
\begin{frame}{Ejemplo 1: velocidad final}\paso{Fórmula:}{v=v_0+at}\paso{Sustitución:}{v=0+(2,5)(8)}\paso{Procedimiento:}{v=20}\res{v=20\,\mathrm{m/s}}\small La velocidad aumenta hacia la derecha.\end{frame}
\begin{frame}{Ejemplo 1: desplazamiento}\paso{Fórmula:}{d=v_0t+\frac{at^2}{2}}\paso{Sustitución:}{d=0(8)+\frac{(2,5)(8)^2}{2}}\paso{Procedimiento:}{d=\frac{160}{2}=80}\res{d=80\,\mathrm{m}}\end{frame}
\begin{frame}{Situación 2: frenado de un vehículo}\centering\auto{1}{-2}{$a=-4\,\mathrm{m/s^2}$}\par\begin{block}{Enunciado}Un automóvil viaja a $24\,\mathrm{m/s}$ hacia el este y frena con aceleración constante de $-4\,\mathrm{m/s^2}$. Determine cuánto tarda en detenerse.\end{block}\datos{$v_0=24\,\mathrm{m/s}$; $v=0$; $a=-4\,\mathrm{m/s^2}$.}\paso{Fórmula y despeje:}{a=\frac{v-v_0}{t}\quad\Longrightarrow\quad t=\frac{v-v_0}{a}}\paso{Sustitución:}{t=\frac{0-24}{-4}}\res{t=6\,\mathrm{s}}\end{frame}
\begin{frame}{Ejemplo 2: distancia de frenado}\datos{$v_0=24\,\mathrm{m/s}$; $v=0$; $a=-4\,\mathrm{m/s^2}$.}\paso{Fórmula:}{v^2=v_0^2+2ad}\paso{Despeje:}{d=\frac{v^2-v_0^2}{2a}}\paso{Sustitución:}{d=\frac{0^2-24^2}{2(-4)}=\frac{-576}{-8}}\res{d=72\,\mathrm{m}}\small El automóvil recorre $72\,\mathrm{m}$ mientras frena.\end{frame}
\begin{frame}{Situación 3: aceleración desconocida}\begin{block}{Enunciado}Una motocicleta aumenta su velocidad de $6\,\mathrm{m/s}$ a $22\,\mathrm{m/s}$ en $4\,\mathrm{s}$. Determine su aceleración y desplazamiento.\end{block}\datos{$v_0=6\,\mathrm{m/s}$; $v=22\,\mathrm{m/s}$; $t=4\,\mathrm{s}$.}\paso{Aceleración:}{a=\frac{22-6}{4}=4\,\mathrm{m/s^2}}\paso{Desplazamiento:}{d=\left(\frac{6+22}{2}\right)(4)=56\,\mathrm{m}}\res{a=4\,\mathrm{m/s^2}\qquad d=56\,\mathrm{m}}\end{frame}
\begin{frame}{Práctica guiada: MRUA horizontal}\begin{enumerate}\item Un vehículo parte del reposo con $a=3\,\mathrm{m/s^2}$ durante $7\,\mathrm{s}$. Calcule $v$ y $d$.\medskip\item Un ciclista pasa de $5\,\mathrm{m/s}$ a $17\,\mathrm{m/s}$ en $6\,\mathrm{s}$. Calcule $a$ y $d$.\medskip\item Un automóvil se desplaza a $30\,\mathrm{m/s}$ y se detiene con $a=-5\,\mathrm{m/s^2}$. Calcule tiempo y distancia de frenado.\end{enumerate}\end{frame}

\section{Movimiento vertical}
\begin{frame}{Movimiento vertical: ¿qué ocurre?}
\begin{columns}[c]
\column{.57\textwidth}
\begin{block}{Idea fundamental}
En ausencia de resistencia del aire, la gravedad produce una aceleración constante dirigida hacia abajo.
\[|g|=9,8\,\mathrm{m/s^2}\]
\end{block}
\begin{itemize}\item Caída libre: se suelta desde el reposo.\item Lanzamiento hacia abajo: comienza con velocidad descendente.\item Lanzamiento hacia arriba: comienza con velocidad ascendente.\end{itemize}
\column{.38\textwidth}
\centering\begin{tikzpicture}[>=Stealth,scale=.85]
\draw[very thick,gray](-.3,0)--(2.5,0);
\filldraw[fill=azul!20,draw=azul](.6,3.5)circle(.23);
\draw[->,ultra thick,naranja](1.2,3.5)--(1.2,1.6)node[right,midway]{$\vec g$};
\draw[dashed,gray](.6,3.2)--(.6,.3);
\node[below]at(1.1,0){Suelo};
\end{tikzpicture}
\end{columns}
\end{frame}
\begin{frame}{Primero: elegir el sentido positivo}
\begin{columns}[T]
\column{.48\textwidth}
\begin{block}{Positivo hacia abajo}
\[a=+9,8\,\mathrm{m/s^2}\]
Útil para objetos que caen o se lanzan hacia abajo.
\end{block}
\column{.48\textwidth}
\begin{block}{Positivo hacia arriba}
\[a=-9,8\,\mathrm{m/s^2}\]
Útil para lanzamientos verticales hacia arriba.
\end{block}
\end{columns}
\vspace{8mm}
\begin{center}\color{morado}\bfseries Las dos convenciones son válidas; no se mezclan en un mismo ejercicio.\end{center}
\end{frame}
\begin{frame}{Formulario: movimiento vertical}
\begin{columns}[T]
\column{.48\textwidth}
\begin{block}{Con el tiempo}
\[v=v_0+at\]
\[h=v_0t+\frac{at^2}{2}\]
\[h=\left(\frac{v_0+v}{2}\right)t\]
\end{block}
\column{.48\textwidth}
\begin{block}{Sin el tiempo}
\[v^2=v_0^2+2ah\]
\[h=\frac{v^2-v_0^2}{2a}\]
\end{block}
\end{columns}
\vspace{4mm}
\centering\small Aquí $h$ representa el desplazamiento vertical con signo según el eje elegido.
\end{frame}
\begin{frame}{Caída libre: reconocer los datos}
\begin{block}{Ejemplo 1}
Se deja caer una piedra desde una altura de $78,4\,\mathrm{m}$. Determine el tiempo de caída. Use el eje positivo hacia abajo.
\end{block}
\etapa{Datos}
\[h=78,4\,\mathrm{m}\qquad v_0=0\,\mathrm{m/s}\]
\[a=+9,8\,\mathrm{m/s^2}\qquad t=?\]
\begin{center}\small\color{verde}«Se deja caer» significa que parte del reposo.\end{center}
\end{frame}
\begin{frame}{Ejemplo 1: tiempo de caída (1/2)}
\etapa{Fórmula}
\[h=v_0t+\frac{at^2}{2}\]
\etapa{Sustitución}
\[78,4=0\cdot t+\frac{9,8t^2}{2}\]
\etapa{Procedimiento}
\[78,4=4,9t^2\]
\end{frame}
\begin{frame}{Ejemplo 1: tiempo de caída (2/2)}
\etapa{Procedimiento: despejar $t$}
\[\frac{78,4}{4,9}=t^2\]
\[16=t^2\]
\[\sqrt{16}=\sqrt{t^2}\]
\[4=t\]
\res{t=4\,\mathrm{s}}
\end{frame}
\begin{frame}{Ejemplo 2: velocidad de impacto}
\etapa{Datos}
\[v_0=0\,\mathrm{m/s},\quad a=+9,8\,\mathrm{m/s^2},\quad t=4\,\mathrm{s}\]
\etapa{Fórmula}
\[v=v_0+at\]
\etapa{Sustitución y procedimiento}
\[v=0+(9,8)(4)\]
\[v=39,2\]
\res{v=+39,2\,\mathrm{m/s}}
\centering\small Positiva: la piedra se mueve hacia abajo, nuestro sentido positivo.
\end{frame}
\begin{frame}{Lanzamiento vertical hacia arriba}
\begin{columns}[c]
\column{.56\textwidth}
\begin{itemize}
\item Elegimos hacia arriba como positivo.
\item Mientras sube: $v>0$.
\item En la altura máxima: $v=0$.
\item Mientras baja: $v<0$.
\item En todo momento: $a=-9,8\,\mathrm{m/s^2}$.
\end{itemize}
\column{.4\textwidth}
\centering\begin{tikzpicture}[>=Stealth,scale=.85]
\draw[thick,gray](-.2,0)--(2.8,0);
\filldraw[fill=azul!20,draw=azul](1,3.6)circle(.2);
\draw[->,ultra thick,verde](.5,.4)--(.5,2.5)node[left,midway]{$v_0$};
\draw[->,ultra thick,naranja](2,3.5)--(2,2)node[right,midway]{$g$};
\node[above]at(1,3.9){$v=0$};
\end{tikzpicture}
\end{columns}
\end{frame}
\begin{frame}{Ejemplo 3: tiempo de subida}
\begin{block}{Enunciado}
Se lanza una pelota verticalmente hacia arriba a $29,4\,\mathrm{m/s}$. Calcule el tiempo de subida.
\end{block}
\etapa{Datos}
\[v_0=+29,4\,\mathrm{m/s}\]
\[v=0\,\mathrm{m/s}\qquad a=-9,8\,\mathrm{m/s^2}\]
\etapa{Fórmula}
\[t=\frac{v-v_0}{a}\]
\end{frame}
\begin{frame}{Ejemplo 3: tiempo de subida, solución}
\etapa{Sustitución}
\[t=\frac{0-29,4}{-9,8}\]
\etapa{Procedimiento}
\[t=\frac{-29,4}{-9,8}\]
\[t=3\]
\res{t_{\mathrm{subida}}=3\,\mathrm{s}}
\end{frame}
\begin{frame}{Ejemplo 4: altura máxima}
\begin{block}{Enunciado}
La pelota del ejemplo anterior sale con $29,4\,\mathrm{m/s}$. Determine la altura máxima alcanzada desde el lanzamiento.
\end{block}
\etapa{Datos}
\[v_0=+29,4\,\mathrm{m/s},\quad v=0\,\mathrm{m/s}\]
\[a=-9,8\,\mathrm{m/s^2},\quad h=?\]
\etapa{Fórmula}
\[h=\frac{v^2-v_0^2}{2a}\]
\end{frame}
\begin{frame}{Ejemplo 4: altura máxima, solución}
\etapa{Sustitución}
\[h=\frac{0^2-(29,4)^2}{2(-9,8)}\]
\etapa{Procedimiento}
\[h=\frac{-864,36}{-19,6}\]
\[h=44,1\]
\res{h_{\max}=44,1\,\mathrm{m}}
\end{frame}
\begin{frame}{Ejemplo 5: tiempo total de vuelo}
\begin{block}{Enunciado}
La pelota sube durante $3\,\mathrm{s}$ y regresa a la misma altura desde donde fue lanzada. Determine el tiempo total.
\end{block}
\etapa{Fórmula}
\[t_{\mathrm{vuelo}}=2t_{\mathrm{subida}}\]
\etapa{Sustitución y procedimiento}
\[t_{\mathrm{vuelo}}=2(3)\]
\res{t_{\mathrm{vuelo}}=6\,\mathrm{s}}
\centering\small Esta relación solo se aplica si regresa a la altura inicial (sin resistencia del aire).
\end{frame}
\begin{frame}{Ejemplo 6: velocidad cuando regresa}
\etapa{Datos}
\[v_0=+29,4\,\mathrm{m/s},\quad a=-9,8\,\mathrm{m/s^2},\quad t=6\,\mathrm{s}\]
\etapa{Fórmula}
\[v=v_0+at\]
\etapa{Sustitución}
\[v=29,4+(-9,8)(6)\]
\etapa{Procedimiento}
\[v=29,4-58,8=-29,4\]
\res{v=-29,4\,\mathrm{m/s}}
\centering\small El signo negativo indica que regresa hacia abajo.
\end{frame}
\begin{frame}{Ejemplo 7: lanzamiento hacia abajo}
\begin{block}{Enunciado}
Una pelota se lanza verticalmente hacia abajo con $v_0=5\,\mathrm{m/s}$. Calcule la velocidad a los $3\,\mathrm{s}$.
\end{block}
\etapa{Datos}
\[v_0=+5\,\mathrm{m/s},\quad a=+9,8\,\mathrm{m/s^2},\quad t=3\,\mathrm{s}\]
\etapa{Fórmula}
\[v=v_0+at\]
\etapa{Sustitución y procedimiento}
\[v=5+(9,8)(3)=5+29,4\]
\res{v=34,4\,\mathrm{m/s}}
\end{frame}
\begin{frame}{Ejemplo 7: desplazamiento hacia abajo}
\etapa{Datos}
\[v_0=5\,\mathrm{m/s},\quad a=+9,8\,\mathrm{m/s^2},\quad t=3\,\mathrm{s}\]
\etapa{Fórmula}
\[h=v_0t+\frac{at^2}{2}\]
\etapa{Sustitución}
\[h=5(3)+\frac{9,8(3)^2}{2}\]
\etapa{Procedimiento}
\[h=15+44,1\]
\res{h=59,1\,\mathrm{m}}
\end{frame}
\begin{frame}{Claves para no equivocarse}
\begin{enumerate}
\item Si se suelta, $v_0=0$.
\item En la altura máxima, $v=0$, pero la gravedad no desaparece.
\item La aceleración siempre apunta hacia abajo.
\item Elegir el eje antes de sustituir los valores.
\item El tiempo total es el doble del de subida \emph{solo si vuelve a la altura inicial}.
\end{enumerate}
\end{frame}
\section{Práctica}
\begin{frame}{Práctica A: MRUA horizontal}\begin{enumerate}\item Un tren pasa de $8\,\mathrm{m/s}$ a $28\,\mathrm{m/s}$ en $10\,\mathrm{s}$. Determine aceleración y desplazamiento.\medskip\item Un automóvil que viaja a $25\,\mathrm{m/s}$ frena uniformemente y se detiene tras $50\,\mathrm{m}$. Calcule aceleración y tiempo de frenado.\medskip\item Una bicicleta parte del reposo y recorre $72\,\mathrm{m}$ en $6\,\mathrm{s}$. Calcule aceleración y velocidad final.\end{enumerate}\end{frame}
\begin{frame}{Práctica B: caída libre}\begin{enumerate}\item Una piedra se suelta desde $45\,\mathrm{m}$. Determine tiempo de caída y rapidez de impacto.\medskip\item Una pelota se lanza hacia abajo a $8\,\mathrm{m/s}$. Determine velocidad y desplazamiento después de $2\,\mathrm{s}$.\medskip\item Un objeto cae desde el reposo durante $5\,\mathrm{s}$. Determine la altura recorrida y la velocidad final.\end{enumerate}\end{frame}
\begin{frame}{Práctica C: lanzamiento vertical}\begin{enumerate}\item Se lanza una pelota hacia arriba a $24,5\,\mathrm{m/s}$. Determine tiempo de subida, altura máxima y tiempo total para regresar al mismo nivel.\medskip\item Un objeto se lanza hacia arriba a $19,6\,\mathrm{m/s}$. Determine su velocidad y desplazamiento después de $3\,\mathrm{s}$.\medskip\item Desde una terraza de $15\,\mathrm{m}$ se lanza una pelota hacia arriba a $12\,\mathrm{m/s}$. Determine el tiempo hasta llegar al suelo.\end{enumerate}\end{frame}
\begin{frame}{Resumen: cómo elegir la fórmula}
\begin{itemize}
\item Si conocemos $v_0$, $a$ y $t$: $v=v_0+at$.
\item Para calcular el desplazamiento con $t$: $h=v_0t+\frac12 at^2$.
\item Si no conocemos $t$: $v^2=v_0^2+2ah$.
\item Si conocemos las dos velocidades y el tiempo: $h=\left(\frac{v_0+v}{2}\right)t$.
\end{itemize}
\vspace{5mm}\centering\color{morado}\bfseries Siempre revisar las unidades y los signos.\end{frame}
\begin{frame}{Cierre de la unidad}\centering\Huge\color{azul}MRUA y movimiento vertical\par\bigskip\normalsize\color{black}Identificar datos $\rightarrow$ elegir eje $\rightarrow$ seleccionar fórmula $\rightarrow$ sustituir $\rightarrow$ interpretar\par\bigskip\color{verde}\large Clases de Matemática Paraíso\end{frame}
\end{document}
